% [原文 p.181]
\chapter{WKB 方法与 Stokes 图}

\ChapterAbstract{谱网络出现在黎曼曲面上线性微分方程的理论中。求解这些微分方程的一种特殊方法是给出形式解的 WKB 方法。利用 Borel 求和，这些形式解可以在一个图的补集上转化为精确解——这个图就是谱网络。我们探讨这一进路，并将它与阿贝尔化过程联系起来。我们将主要限于 $\SL_2(\C)$ 的情形。}

\section{概述}

本章中，我们将从微分方程的分析出发，看到关于平坦联络的阿贝尔化过程以及谱网络之出现的另一个视角。

大致的推理路线如下：
\[
\boxed{\begin{array}{c}\text{秩 }n\text{ 丛上的}\\[-2pt]\text{平坦联络}\end{array}}
\longrightarrow
\boxed{\begin{array}{c}\text{平坦截面的}\\[-2pt]\text{微分方程}\end{array}}
\xrightarrow{\ \text{WKB 方法}\ }
\boxed{\text{阿贝尔化}}
\]

由 Riemann–Hilbert 对应（见第 3.4 节），特征簇描述平坦联络的规范等价类。黎曼曲面 $S$ 上秩为 $n$ 的丛上的 $\SL_n(\C)$-联络是平坦的，当且仅当它容许由平坦截面组成的局部基。平坦截面的方程导出一个在 $S$ 上有整体意义的微分方程。把 WKB 方法应用于这个微分方程，自然地引出带线丛与阿贝尔联络的谱曲线。于是我们恢复了上一章第 15 章中的阿贝尔化过程。

从物理的观点看，WKB 方法在经典物理与量子物理之间建立联系。更确切地说，它把量子力学中的波函数（即某个微分方程——典型的是

% [原文 p.182]
Schrödinger 方程——的解）与经典相空间某个拉格朗日子流形的生成函数联系起来。这个过程是可逆的：从经典数据出发，我们可以恢复波函数。

WKB 方法可以形式地进行，也可以精确地进行。形式 WKB 方法把波函数重构为形式幂级数，而精确 WKB 方法给出良定义的解析函数。所涉及的主要概念概括于下图：
\[
\begin{array}{cc}
\multicolumn{2}{c}{\boxed{\text{WKB 方法}}}\\[7pt]
\text{形式}\swarrow & \searrow\text{精确}\\[7pt]
\boxed{\begin{array}{c}\text{Hamilton--Jacobi}\\[-2pt]\text{方程}\end{array}} & \boxed{\begin{array}{c}\text{Borel 求和}\\[-2pt]\text{方法}\end{array}}\\[7pt]
\downarrow & \downarrow\\[7pt]
\boxed{\text{递推}} & \boxed{\begin{array}{c}\text{Stokes 图 =}\\[-2pt]\text{谱网络}\end{array}}
\end{array}
\]

WKB 方法给出的形式解，来自求解一个 Hamilton–Jacobi 方程（在黎曼曲面的情形这只是一个多项式方程）与一个递推。这些解通常收敛半径为 0。为求得精确解，要使用 Borel 求和。对 $\SL_2(\C)$，Borel 求和方法在 $S\setminus\Gamma$ 上收敛，其中 $\Gamma$ 是一个特殊的图，称为 Stokes 图；它是 $A_1$ 型谱网络。猜想这对更高秩仍然成立。

为保持在本书范围内，我们只能勾画主要思想，而不给出技术性很强的证明。Borel 可和性与精确 WKB 这一研究领域非常活跃，许多猜想仍未解决，尤其是关于阶至少为 3 的微分方程。

我们邀请感兴趣的读者参阅以下文献以了解细节。这里只给出很不完整的成果清单；更完整的文献综述可在 [289] 中找到。

最初的 WKB 方法旨在给出量子力学中方程的近似解，由 Wentzel [357]、Kramers [241] 与 Brillouin [68] 于 1926 年各自独立发展。1980 年代初，Voros [352] 引入了精确 WKB 方法。大约同时，Ecalle [123] 引入了可复苏函数（resurgent function）空间，并发展了所谓的外星演算（alien calculus）。关于 Ecalle 的复苏理论，我们还推荐 Sauzin [310] 的漂亮导引。对物理背景更强的读者，我们推荐 Dorigoni [121] 与 Serone [325] 的讲义，以及 Mariño 的书 [258, Chap. 2]。较近期的发展包括与谱网络、阿贝尔化以及簇 Fock–Goncharov 坐标的联系。特别地，精确 WKB 方法与

% [原文 p.183]
谱网络之间的关系最早由 Hollands–Neitzke [207] 详细给出；我们还推荐 Gaiotto–Moore–Neitzke 的基础性论文 [154]、Hollands–Rüter–Szabo [208, Sect. 4.4]、Iwaki–Nakanishi 的工作 [216]、Kawai–Takei 的书 [229]，以及 Nikolaev 近期的几何观点 [290]。

\section{形式局部 WKB}

WKB 方法以 G. Wentzel、H. Kramers 与 L. Brillouin 命名，1926 年发展起来，用以求 Schrödinger 方程的近似解。更一般地，它提供了求解任何带参数 $h$（视作小参数）的线性微分方程的工具；这类方程通常来自量子力学，其中 $h$ 是 Planck 常数。

考虑一个线性微分方程。历史上的例子是时间无关（或定态）的一维 Schrödinger 方程：
\begin{equation}
(h^2\partial^2-\hat t(z,h))\psi(z,h)=0.
\tag{16.1}
\end{equation}
这里 $z$ 是 $\C$ 的一个开集上的局部坐标，$\partial=\frac{\mathrm{d}}{\mathrm{d}z}$ 是导数，$h$ 是（小）参数，$\hat t(z,h)$ 是给定函数（$h$ 的 Taylor 级数），而 $\psi$ 是我们要解的函数。

更一般地，我们可以考虑 $n$ 个变量 $x_1,x_2,\dots,x_n$、导数 $\partial_i:=\frac{\partial}{\partial x_i}$ 的任意线性微分方程。这样的方程可写成如下紧凑形式：
\begin{equation}
P(h\partial_1,\dots,h\partial_n)\psi(x_1,\dots,x_n,h)=0,
\tag{16.2}
\end{equation}
其中 $P$ 是在算子 $h\partial_i$ 处取值的多项式。这里 $h\in\C^*$ 仍是（小）参数，$P$ 的系数是 $x_i$ 与 $h$ 的函数。在形式设定中，我们假设系数是 $h$ 的形式幂级数。把 $P$ 在偏导数处取值，意思是把变量替换为导数，同时把 $P$ 的系数放在左边。

对 Schrödinger 方程 (16.1)，有 $n=1$、$x_1=z$、$P(X)=X^2-\hat t(z,h)$。确实，把 $X$ 换成 $h\partial$ 就得到 (16.1)。下文中我们把 $\psi(x_i,h)$ 用作 $\psi(x_1,\dots,x_n,h)$ 的简写。

WKB 方法研究如下拟设，以得到形式解：
\begin{equation}
\psi(x_i,h)=\exp\left(\frac1h\left(s_0(x_i)+hs_1(x_i)+h^2s_2(x_i)+\dots\right)\right).
\tag{16.3}
\end{equation}
我们在 Schrödinger 方程上说明这一方法。

\begin{example}[16.1] 再次考虑一个变量 $x_1=z$ 的 Schrödinger 方程：
\begin{equation}
(h^2\partial^2-\hat t(z,h))\psi(z,h)=0,
\tag{16.4}
\end{equation}
% [原文 p.184]
其中 $z\in U$，$U$ 是 $\C$ 的开集，$\partial=\frac{\partial}{\partial z}$，而 $\hat t(z,h)$ 是系数为 $U$ 上全纯函数的形式幂级数。代入拟设 $\psi=\exp\left(\frac1h s(z)\right)$，其中 $s(z)=s_0(z)+hs_1(z)+\dots$，得到
\begin{equation}
(\partial s)^2+h\partial^2s-\hat t(z,h)=0.
\tag{16.5}
\end{equation}
写 $\hat t=t+O(h)$，则 (16.5) 的 $h^0$-项给出
\[
(\partial s_0)^2-t=0,
\]
由此得到两个解 $\partial s_0=\pm\sqrt{t}$。WKB 方法的重要特点是：只要知道 $\hat t(z,h)=\sum_i h^it^{(i)}(z)$ 的所有阶，就能从 $s_0$ 递归地得到所有 $s_i$（$i>0$）。例如，(16.5) 的 $h$-项读作 $2\partial s_0\partial s_1+\partial^2s_0-t^{(1)}=0$。一般而言，$s_i$ 的方程是
\[
2\partial s_0\partial s_i+\sum_{k=1}^{i-1}\partial s_k\partial s_{i-k}+\partial^2s_{i-1}-t^{(i)}=0.
\]
最终我们得到两个基本形式解，对应于 $\partial s_0$ 的每个选择；这与 Schrödinger 方程是二阶方程这一事实相符。$s_i$ 的积分常数可被吸收进一个整体因子（$h$ 的形式幂级数）中，该因子给基本解乘上标量。\end{example}

% [Tutor Note: 把 $\psi=\exp(s/h)$ 代入 $(h^2\partial^2-\hat t)\psi=0$ 复核 (16.5)：$\partial\psi=(s'/h)\psi$，$\partial^2\psi=(s''/h+s'^2/h^2)\psi$，故 $h^2\partial^2\psi=(hs''+s'^2)\psi$，方程化为 $s'^2+hs''=\hat t$，正是 (16.5)。再逐阶核对：$h^0$ 阶给出 $(s_0')^2=t^{(0)}$；$h^i$ 阶（$i\ge1$）给出 $2s_0's_i'+\sum_{k=1}^{i-1}s_k's_{i-k}'+s_{i-1}''=t^{(i)}$，即正文中对 $s_i$ 的递推式，其中交叉项来自 $(s')^2$ 的展开、二阶导数项来自 $h s''$。]

一般地，把 WKB 拟设 (16.3) 代入微分方程 (16.2)，最低阶项给出关于 $s_0$ 的非线性微分方程。所有 $i>0$ 的 $s_i$ 随后由 $s_0$ 递归确定。$s_0$ 的方程为
\begin{equation}
P(x_i,\partial_is_0)=0.
\tag{16.6}
\end{equation}
这就是所谓的 Hamilton–Jacobi 方程。该方程在辛几何中有漂亮的解释，给出一般的求解方案 [32, Chap. 9]。

\section{插曲：Hamilton–Jacobi 方程}

我们描述如何把 Hamilton–Jacobi 方程化归为一对彼此相关的方程，一个是多项式方程，另一个是常微分方程。我们假设读者具备一些辛几何知识。由于此后我们只在单变量下工作，本节对后文并非必需。更多细节，我们推荐 Arnold 的书 [32]，尤其是其第 9 章。

考虑 $2n$ 维辛流形 $(M,\omega)$。局部地可取 Darboux 坐标 $(x_i,y_i)_{1\le i\le n}$。设 $K\subset M$ 是由方程 $P(x_i,y_i)=0$ 局部给出的超曲面。目标是找出 $K$ 内的一个拉格朗日子流形 $L$。这等价于求解 Hamilton–Jacobi 方程。确实，局部地任何拉格朗日子流形都可通过只依赖一半变量的生成函数 $s(x_i)$ 描述。拉格朗日流形 $L$ 然后由 $y_i=\partial_is$ 定义。$L\subset K$ 的

% [原文 p.185]
条件于是变为
\begin{equation}
P(x_1,\dots,x_n,\partial_1s,\dots,\partial_ns)=0,
\tag{16.7}
\end{equation}
它恰恰就是 Hamilton–Jacobi 方程 (16.6)。

我们可以用另一种方式构造 $L$，从而得到 Hamilton–Jacobi 方程的解。首先须确保微分方程的初始条件落在 $K$ 中，这归结为一个多项式方程。其次对 $P$ 的辛梯度流作积分，这等价于解一个常微分方程。结果就是拉格朗日流形 $L$。

使用适当的坐标，(16.7) 的初始条件可描述为 $s|_{\{x_1=0\}}=f(x_2,\dots,x_n)$，其中 $f$ 是某个给定函数。我们必须确保
\[
P(0,x_2,\dots,x_n,\partial_1s,\partial_2f,\dots,\partial_nf)=0,
\]
这是关于 $\partial_1s|_{\{x_1=0\}}$ 的多项式方程。后者的解确定 $K$ 中一个 $(n-1)$ 维的各向同性子流形 $L'$。

然后考虑由 $P$ 的辛梯度生成的流。物理上，这对应于以 $P$ 为 Hamilton 量的经典系统的时间演化。由于 $K$ 是 $P$ 的一个等值集，该流保持在 $K$ 内。对这一流作积分归结为解一个常微分方程。$L'$ 在该流下的像是一个拉格朗日子流形 $L$。确实，各向同性子空间在 Hamilton 流下的像仍各向同性，而 $P$ 的流方向是 $\omega|_K$ 的核，因为 $K$ 是 $P$ 的等值集。因此 $\omega|_L=0$ 且 $\dim(L)=n$，故 $L$ 是拉格朗日的。

\section{形式整体 WKB}

考虑闭黎曼曲面 $S$。我们要在 $S$ 上理解微分方程 (16.2) 的含义，然后应用形式 WKB 方法。

由于现在只处理单个（复）变量 $z$，设定大大简化。局部地 (16.2) 现在读作
\begin{equation}
(h^n\partial^n+a_1(z,h)h^{n-1}\partial^{n-1}+\dots+a_n(z,h))\psi(z,h)=0,
\tag{16.8}
\end{equation}
其中 $\partial=\frac{\mathrm{d}}{\mathrm{d}z}$，$a_i$ 是 $z$ 的全纯函数、$h$ 的形式幂级数，而 $\psi$ 是我们要解的未知函数。

为赋予该微分方程整体意义，我们先把其改写为一阶微分方程：
\begin{equation}
(h\partial-A)\Psi=0,
\tag{16.9}
\end{equation}
其中

% [原文 p.186]
\begin{equation}
A=\begin{pmatrix}0&1&&\\&\ddots&\ddots&\\&&0&1\\-a_n&-a_{n-1}&\cdots&-a_1\end{pmatrix}
\qquad\text{和}\qquad
\Psi=\begin{pmatrix}\psi<本地路径>
\tag{16.10}
\end{equation}
我们把 $h\partial-A$ 解释为 $S$ 上某个全纯丛中的 $h$-联络。

回忆第 3.2 节：联络是定义截面方向导数的一种方式，即它把截面 $s$ 与向量场 $X\in\Gamma(TS)$ 联系为一个新截面 $\nabla_X(s)$。$h$-联络非常类似，唯一差别在于修改过的 Leibniz 法则：

\begin{definition}[16.2] $S$ 上向量丛 $E$ 的一个 $h$-联络是映射 $\nabla:\Gamma(E)\times\Gamma(TS)\to\Gamma(E)$，$(s,v)\mapsto\nabla_v(s)$，满足
\begin{enumerate}
  \item 纤维中的线性性：$\nabla_v(s+t)=\nabla_v(s)+\nabla_v(t)$；
  \item 对 $v$ 的 $C^\infty(S)$-线性性：$\nabla_{fv+gw}=f\nabla_v+g\nabla_w$；
  \item $h$-Leibniz 法则：$\nabla_v(fs)=h\,\mathrm{d}f_v(s)+f\nabla_v(s)$，
\end{enumerate}
其中 $s,t\in\Gamma(E)$，$f,g\in C^\infty(S)$，$v,w\in\Gamma(TS)$。\end{definition}

当 $h=1$ 时，我们恢复通常的联络概念。

\begin{example}[16.3] 再次考虑 Schrödinger 方程
\[
(h^2\partial^2+t(h,z))\psi(h,z)=0.
\]
把它改写为一阶微分方程
\[
\left(h\partial-\begin{pmatrix}0&1\\-t&0\end{pmatrix}\right)\begin{pmatrix}\psi<本地路径>
\]
该算子是全纯丛 $E$ 上的一个 $h$-联络。由于我们希望微分方程有两个线性无关的解，该联络必须是平坦的。这蕴含 $\bar\partial t=0$，从而解释了为何考虑全纯联络。此外，丛 $E$ 是平凡丛 $\OO$，因为我们可以用平坦截面的基把它平凡化。

让我们确定 $\psi$ 的整体意义。它是某个线丛 $L$ 的截面。于是 $h\partial\psi$ 是 $L\otimes K$ 的截面，其中 $K$ 表示全纯 1-形式的典范丛。因此 $\Psi$ 是 $L\oplus L\otimes K=E\simeq\OO$ 的截面。相应的行列式丛也是平凡的：$L^2\otimes K\simeq\OO$。故 $L=K^{-1/2}$，其中我们必须选取 $K$ 的一个平方根（即一个线丛 $\ell$，使得 $\ell\otimes\ell\simeq K$）。\end{example}

上面例子中的论证可以推广。方程 (16.9) 的整体意义如下：$h\partial-A$ 是秩为 $n$ 的平凡全纯丛 $\OO$ 中的全纯 $h$-联络，而我们要找平坦的全纯截面 $\Psi$。

% [原文 p.187]
初始的未定函数 $\psi$ 是线丛 $K^{(1-n)/2}$ 的全纯截面。确实，它是某个线丛 $L$ 的截面。于是 $h^k\partial^k\psi$ 是 $L\otimes K^k$ 的截面，所以 $\Psi$ 是 $\bigoplus_{k=0}^{n-1}L\otimes K^k$ 的截面。相应的行列式丛由 $L^n\otimes K^{(n-1)n/2}$ 给出，它是平凡丛 $\OO$。故 $L=K^{(1-n)/2}$。

我们在 (16.3) 中看到，WKB 方法寻找如下形式的形式局部解
\begin{equation}
\psi(h,z)=\exp\left(\frac1h s(z,h)\right),
\tag{16.11}
\end{equation}
其中 $s(z,h)=s_0(z)+hs_1(z)+h^2s_2(z)+\dots$ 是形式幂级数。下面我们给出这一拟设的整体意义。

最低阶项给出 Hamilton–Jacobi 方程 $P(s_0,z)=0$，在此情形它是多项式方程。这一方程应理解为定义了一个分歧覆盖 $\pi:\Sigma\to S$，称为谱曲线，$s_0$ 在其上有良好定义：
\begin{equation}
\Sigma=\{(p,z)\in T^*S\mid P(p,z)=0\},
\tag{16.12}
\end{equation}
其中 $p$ 表示 $T^*S$ 纤维上的全纯坐标。

谱曲线 $\Sigma$ 是次数为 $n$ 的分歧覆盖，$n$ 是多项式 $P$（以及初始线性微分方程）的次数。局部地在分歧点之外，可把各叶标号为 $1$ 到 $n$。利用割线可以定义整体标号。

WKB 拟设的整体意义如下。我们寻找 $n$ 个不同形式解
\begin{equation}
\psi^{(i)}_{\mathrm{formal}}(z,h)=\exp\left(\frac1h\int_{y_0}^{y_i}\tilde\lambda(h)\right),
\tag{16.13}
\end{equation}
其中 $y_i\in\pi^{-1}(z)$ 是 $z$ 在 $\Sigma$ 第 $i$ 叶中的原像，$y_0\in\Sigma$ 是基点，而 $\tilde\lambda(h)=\lambda_0+h\lambda_1+\dots$ 是 1-形式的形式幂级数。

重要的观察是：$\lambda_0$ 是 $T^*S$ 的 Liouville 1-形式限制到谱曲线 $\Sigma$ 上的结果。换言之：仅用 (16.13) 中的 Liouville 1-形式，就把初始微分方程解到最高阶。其他 1-形式 $\lambda_i$（$i\ge1$）以递归方式确定。

\begin{example}[16.4] 对 Schrödinger 方程
\[
(h^2\partial^2-t(z,h))\psi(h,z)=0,
\]
其中 $t(z,h)=t_0(z)+ht_1(z)+\dots$，WKB 方法 (16.11) 给出
\[
s_0^2(z)-t_0(z)=0.
\]
谱曲线 $\Sigma$ 由 $p^2=t_0$ 定义。局部地选取平方根 $\sqrt{t_0(z)}$ 的一个分支，等同于选取 $\Sigma\to S$ 的一个分支。

限制到 $\Sigma$ 上的 Liouville 1-形式于是为 $\lambda_0=\sqrt{t_0}$。按 (16.13)，可直接验证

% [原文 p.188]
\[
\exp\left(h^{-1}\int^z\sqrt{t_0(z)}\right)
\]
把 Schrödinger 方程解到首阶。$\tilde\lambda$ 的次领头项由
\[
\tilde\lambda=\sqrt{t_0}-h\frac{\partial t_0}{4t_0}-h^2\sqrt{t_0}\frac{5(\partial t_0)^2-4t_0\partial^2t_0}{32t_0^3}+\dots
\]
给出。\end{example}

一般地，容易验证形式解 $\psi^{(i)}_{\mathrm{formal}}$ 是
\begin{equation}
(h\partial-\lambda^{(i)})\psi^{(i)}=0
\tag{16.14}
\end{equation}
的解，其中 $\lambda^{(i)}$ 是 $\tilde\lambda$ 限制到 $\Sigma$ 第 $i$ 叶的结果。我们可以这样解释该方程：$h\partial-\lambda^{(i)}$ 是谱曲线 $\Sigma$ 第 $i$ 叶的一个开集 $U$ 上某个线丛 $L$ 上的 $h$-联络，而 $\psi^{(i)}$ 是 $L$ 的截面。于是我们看到，(16.14) 是 (16.9) 中平坦联络 $h\partial-A$ 的一个阿贝尔化。

WKB 拟设 (16.12) 在 $S\setminus P$ 上给出 $n$ 个形式解，其中 $P$ 表示极点集与分支点集。这里似乎有一个矛盾：在分支点附近各叶会互换，但从微分方程的角度看，这些是光滑点，不应有单值性。这一表面矛盾的解决在于：WKB 解只是收敛半径为零的形式解。

天真的希望是：存在微分方程的解析解 $\psi^{(i)}_{\mathrm{an}}$，使得当 $h\to0$ 时 WKB 形式解 $\psi^{(i)}_{\mathrm{formal}}$ 是它的渐近展开。不幸的是，这样的精确解一般并不存在。我们必须限制 $h$ 趋于 0 的方式。由于 WKB 拟设中的指数函数，我们必须限制 $h$ 停留在某个半平面内，如下例所示。

\begin{example}[16.5] 考虑函数 $f(z)=\ee^{-1/z}$。它是 $\C^*$ 上的全纯函数，在原点有本性奇点。当 $\operatorname{Re}(z)>0$ 时，$f(z)$ 在 $z\to0$（同时停留在 $\operatorname{Re}(z)>0$ 所定义的半平面内）时的极限为 0。当 $\operatorname{Re}(z)<0$ 时，同一极限发散。\end{example}

下文中，我们将把 WKB 形式解修改为精确解，这些精确解一般在某个半平面内渐近于形式解。曲面上这种渐近行为失效的点构成一个组合对象——谱网络。

\section{通过 Borel 求和的精确 WKB}

关于 WKB 方法的一个重要问题是收敛性。一般而言，(16.3) 中的形式和从不收敛；它只是一个渐近级数。1980 年代，Voros [352] 引入了精确 WKB 方法，它把形式 WKB 方法与 Borel 求和结合起来以得到有效的解。在

% [原文 p.189]
介绍 Laplace 变换的一些基础之后，我们按 [216] 阐述 Borel 求和方法。物理学中有大量例子，这些方法在那里以复苏（resurgence）或外星演算之名出现，例如见 [257]，或 [121] 中更数学化的处理。

函数 $f:\R^+\to\C$ 的 Laplace 变换由
\begin{equation}
\mathcal{L}[f](s)=\int_0^\infty f(t)\ee^{-st}\mathrm{d}t
\tag{16.15}
\end{equation}
给出，其中 $s\in\C$。对 $\operatorname{Re}(s)>0$ 以及满足 $|f(t)|\le C_1\ee^{C_2t}$（对某些正常数 $C_1,C_2$）的函数 $f$，这是良定义的。

可以把变换 (16.15) 看成 Fourier 变换的解析延拓；取 $s=\ii\xi$ 时我们几乎\footnote{Fourier 变换是 $s=\ii\xi$ 处的双边 Laplace 变换，即一种在 $\R$ 上积分的 Laplace 变换。}得到 Fourier 变换。

\begin{example}[16.6] 直接计算表明，$f(t)=t^q$（$q\in\C$，$\operatorname{Re}(q)>-1$）的 Laplace 变换由
\[
\mathcal{L}[t^q](s)=\frac{\Gamma(q+1)}{s^{q+1}}
\]
给出，其中 $\Gamma$ 是 Gamma 函数。特别地，对 $q=n\in\N$：
\begin{equation}
\mathcal{L}[t^n](s)=\frac{n!}{s^{n+1}}.
\tag{16.16}
\end{equation}
\end{example}

与 Fourier 变换一样，Laplace 变换对求解微分方程非常有用。用分部积分可见
\[
\mathcal{L}[f'](s)=s\mathcal{L}f(s)-f(0).
\]
因此把 $\mathcal{L}$ 作用于线性微分方程就把它变为多项式方程。

此外，Laplace 变换可通过围道积分求逆：
\[
(\mathcal{L}^{-1}f)(t)=\int_\gamma \ee^{st}f(s)\mathrm{d}s.
\]
积分路径 $\gamma$ 通常形如 $\{z\in\C\mid\arg(z)=\theta_0\}$（对某个固定的 $\theta_0\in[0,2\pi)$），且必须落在被积函数的收敛区域内。

逆 Laplace 变换可用于求解常系数非齐次线性微分方程。为求解
\begin{equation}
Df(t)=g(t),
\tag{16.17}
\end{equation}

% [原文 p.190]
其中 $D$ 是常系数线性微分算子，只需找到 Green 函数 $G$，即 $DG(t)=\delta_0(t)$ 的解，其中 $\delta_0$ 是支撑在零点的 Dirac 分布。于是 (16.17) 的通解是卷积积 $G\star g$。把 Laplace 变换作用于 $DG(t)=\delta(t)$ 给出 $P(t)\mathcal{L}(G)=1$，其中 $P$ 是满足 $D=P\left(\frac{\mathrm{d}}{\mathrm{d}t}\right)$ 的多项式。因此
\[
G(t)=\mathcal{L}^{-1}(1/P).
\]
当 $P$ 有单零点时，$1/P$ 有 $n$ 个极点，于是逆 Laplace 变换有 $n+1$ 条不同（同伦型的）积分路径。这给出 $n$ 个线性无关的 Green 函数，因为所有留数之和为零。注意 Fourier 变换只给出一个 Green 函数。由逆 Laplace 变换得到的 Green 函数相差 $\ee^{st}/P(s)$ 的留数。这些留数恰是齐次方程 $Df(t)=0$ 的基本解。这样的留数是 $\ee^{x_0t}$ 的常数倍，其中 $x_0$ 是 $P$ 的一个根。于是我们恢复了常系数齐次微分方程的求解方案。

\begin{example}[16.7] 考虑 $f''(t)+f(t)=\delta_0(t)$。施加 Laplace 变换并令 $F=\mathcal{L}f$ 得 $(s^2+1)F=1$。$\frac{\ee^{st}}{1+s^2}$ 在 $\pm\ii$ 处的留数由 $\pm\frac{\ee^{\pm\ii t}}{2\ii}$ 给出，它们构成齐次方程 $f''(t)+f(t)=0$ 的解的一组基。\end{example}

现在我们可以解释 Borel 求和，它是一个两步过程。第一步是 Borel 变换。

\begin{definition}[16.8] 形式幂级数 $f(s)=\sum_{n\ge0}f_ns^{-n}$ 是 Borel 可和的，如果
\begin{equation}
f_B(y)=\sum_{n\ge1}f_n\frac{y^{n-1}}{(n-1)!}
\tag{16.18}
\end{equation}
在 $y=0$ 附近收敛，并且可以解析延拓到包含正实轴的一个区域 $\Omega$，在其中有 $|f_B(y)|\le C_1\ee^{C_2|y|}$（对某些正常数 $C_1,C_2$）。映射 $f\mapsto f_B$ 称为 Borel 变换。\end{definition}

把定义 (16.18) 与 (16.16) 中函数 $t\mapsto t^n$ 的 Laplace 变换相比较，我们看出 Borel 变换为何是逐项形式上的逆 Laplace 变换。

Borel 可和函数的一个典型例子是形式幂级数 $\sum_{n\ge0}f_ns^{-n}$，其系数满足 $|f_n|\le Cn!\rho^n$（对某些正常数 $C$ 与 $\rho$）。这样的形式幂级数称为 Gevrey-1 的。对 Gevrey-1 幂级数，Borel 变换的收敛半径严格为正。

重要的是，Borel 可和函数在物理的微扰方法中无处不在。一个典型例子是 Feynman 图展开，因为已知含 $n$ 个顶点的 Feynman 图约有 $n!$ 个。

Borel 求和过程的第二步，是把 Laplace 变换作用于第一步（解析延拓后的）Borel 变换：

% [原文 p.191]
\begin{definition}[16.9] Borel 可和函数 $f(s)=\sum_{n\ge0}f_ns^{-n}$ 的 Borel 和，或 Borel 求和，是由
\begin{equation}
S[f](s)=f_0+\int_0^\infty\ee^{-sy}f_B(y)\mathrm{d}y
\tag{16.19}
\end{equation}
给出的 $s$ 的函数。\end{definition}

注意该定义对 $\operatorname{Re}(s)>0$ 良定义，因为 $\R_+$ 包含在 $f_B$ 的定义域中，且该定义域中有指数界。

\begin{example}[16.10] 考虑 $f(s)=\sum_{n\ge0}s^{-n}$，它在 $|s^{-1}|<1$ 时收敛。Borel 变换给出 $f_B(y)=\sum_{n\ge1}\frac{y^{n-1}}{(n-1)!}=\ee^y$。于是 Borel 和由 $S[f](s)=1+\int_0^\infty\ee^{y(1-s)}\mathrm{d}y=1+\frac{1}{1-s}=\frac{1}{1-s^{-1}}$ 给出，其中 $\operatorname{Re}(s^{-1})<1$。注意 $S[f]$ 是 $f$ 的解析延拓。\end{example}

% [Tutor Note: 例 16.10 的中间项原文印作 $1+\frac{1}{1-s}$，但与两端不相容：$\int_0^\infty\ee^{y(1-s)}\mathrm{d}y=\frac{1}{s-1}$（在 $\operatorname{Re}s>1$ 时收敛），故应有 $1+\frac{1}{s-1}=\frac{s}{s-1}=\frac{1}{1-s^{-1}}$；而 $1+\frac{1}{1-s}=\frac{2-s}{1-s}\ne\frac{1}{1-s^{-1}}$。故中间分式的分母应为 $s-1$，原文是一处符号笔误，照译并在此指出。]

\begin{example}[16.11] 考虑微分方程
\[
-\frac{\mathrm{d}f}{\mathrm{d}s}+f(s)=1/s,
\]
即所谓的一阶 Euler 微分方程。一个形式解由幂级数 $f(s)=\sum_{n\ge0}(-1)^nn!s^{-n-1}$ 给出。容易验证 $f$ 是 Borel 可和的，且
\[
S[f](s)=\int_0^\infty\frac{\ee^{-su}}{1+u}\mathrm{d}u,
\]
它给出 Euler 方程的精确解。\end{example}

% [Tutor Note: 可核对例 16.11：设 $f(s)=\sum_{n\ge0}(-1)^nn!s^{-n-1}$，则 $-f'+f$ 中 $s^{-m-1}$（$m\ge1$）的系数为 $(-1)^{m-1}m!+(-1)^mm!=0$，只剩 $m=0$ 的项 $s^{-1}$，故 $-f'+f=1/s$，即 $f$ 确是 Euler 方程的形式解。再核对 Borel 变换：把级数按 $f(s)=\sum_{m\ge1}f_ms^{-m}$ 重排得 $f_m=(-1)^{m-1}(m-1)!$，故按 (16.18) 有 $f_B(y)=\sum_{m\ge1}(-1)^{m-1}y^{m-1}=\frac{1}{1+y}$，从而 $S[f](s)=\int_0^\infty\frac{\ee^{-su}}{1+u}\mathrm{d}u$，与原文一致。]

Borel 求和的一个重要特点是：对已经收敛的级数，它是恒等映射：

\begin{proposition}[16.12] 若 $f(s)$ 在 $s=\infty$ 附近收敛到全纯函数，则 Borel 求和是恒等映射：$S[f]=f$。\end{proposition}

证明只需看出：若 $f$ 在 $s=\infty$ 附近收敛，则第一步确实是逆 Laplace 变换，而它被第二步抵消。

在函数 $f$ 不收敛的情形，它给出 Borel 求和的渐近级数：

\begin{proposition}[16.13（{{[216, Prop. 2.11(b)]}}）] 当 $s\to\infty$ 时，$S[f](s)\sim f(s)$。\end{proposition}

最后，直接由定义得到的若干性质如下：

\begin{proposition}[16.14（{{[216, Prop. 2.11(a, c)]}}）]
\begin{enumerate}
  \item 对 Borel 可和的 $f$ 与 $g$，$S[f+g]=S[f]+S[g]$，且 $S[fg]=S[f]S[g]$；
  \item 若 $A(t)=\sum_{k\ge0}A_kt^k$ 在 $t=0$ 附近收敛，且 $f(s)=\sum_{n\ge1}f_ns^{-n}$ 无常数项，则 $S[A\circ f]=A\circ S[f]$。
\end{enumerate}
\end{proposition}

\begin{remark}[16.15] 有一类函数称为可复苏的（resurgent），我们在上面用到的运算在其中封闭：Laplace 变换、Borel 变换、复合，以及在解析延拓与卷积下封闭。粗略地说，可复苏函数是 Gevrey-1 的，且其 Borel 变换沿从原点出发的射线上只有有限多个奇点。更多细节见 [121, 123, 310]。\end{remark}

% [原文 p.192]
Borel 求和有许多变化与推广。例如，我们可以对形式幂级数之外更一般的函数定义求和。对形如 $g(s)=\ee^{s\mu}s^{-\rho}f(s)$ 的形式函数 $g$（$\mu,\rho\in\C$，$f$ 同上），我们定义
\begin{equation}
S[g](s):=\ee^{s\mu}s^{-\rho}S[f](s).
\tag{16.20}
\end{equation}
另一个推广是在第二步中加入某个角度。不必在 $\R_+$ 上积分，我们可以定义方向 $\theta$ 上的 Borel 和，其中 $\theta\in[0,2\pi)$：
\[
S_\theta[f](s):=f_0+\int_0^{\infty\ee^{-\ii\theta}}\ee^{-sy}f_B(y)\mathrm{d}y,
\]
其中我们在从原点出发、与实轴成角 $-\theta$ 的半直线上积分。关于 $f$ 在方向 $\theta$ 上 Borel 可和，有相应的概念。Borel 和 $S_\theta[f]$ 在 $\{s\in\C\mid|\arg(s)-\theta|<\pi/2,\ |s|\gg1\}$ 上解析。

当 $f_B$ 在积分路径上有极点时，带角 $\theta$ 的 Borel 求和尤为重要。此时有两种选择：从极点的上方或下方绕过（见图 16.1）。两个结果相差绕极点的留数。这就是所谓的 Stokes 现象。

\Fig[0.75\linewidth]{figures/fig_16_1.png}{16.1}{Borel 求和的围道（红点表示极点）。}

Stokes 现象来自具有给定渐近级数的解析函数的不唯一性。这表现为：Borel 求和依赖于积分路径（第二步 Laplace 变换中的路径）。两条非同伦的路径一般给出不同结果。

\begin{example}[16.16] 考虑微分方程
\[
f'(s)+f(s)=\frac1s.
\]
一个形式解由 $f(s)=\sum_{n\ge0}n!s^{-n-1}$ 给出，它是 Borel 可和的。Borel 变换为 $f_B(y)=\frac{1}{1-y}$，在实轴上 $y=1$ 处有极点。可以算出，从极点上方的 Borel 和 $S_{0^+}[f]$ 与从极点下方的 Borel 和 $S_{0^-}[f]$ 之差由

% [原文 p.193]
\[
S_{0^+}[f](s)-S_{0^-}[f](s)=2\pi\ii\ee^{-s}
\]
给出。注意该差是指数小的。\end{example}

我们可以把上例稍作推广。假设 Borel 变换 $f_B(y)$ 在 $y=y_0\in\R$ 处有极点。则函数 $\ee^{-sy}f_B(y)$ 在 $y=y_0$ 处的留数由 $\ee^{-sy_0}\operatorname{Res}_{y_0}(f_B)$ 给出，它在 $s\to\infty$ 时是指数小的。在 Borel 和中采用从 $y_0$ 上方或下方经过的积分路径，得到两个不同的解，它们在 $s\to\infty$ 时有相同的渐近，因为二者之差是指数小的。

\begin{remark}[16.17] Stokes 现象可由所谓的外星导数（alien derivative）捕捉，例如见 [121, Sects. 4 and 6]。它是与描述跨越 Stokes 线之跳跃的自同构相联系的导子。外星导数可通过 Ecalle 的桥方程与通常的导数联系起来。

复苏理论可视为对 Stokes 现象的系统研究。它可以赋予不可 Borel 求和的渐近函数以意义。关键思想是寻找如下形式的微分方程的解
\[
\tilde F(h)=\sum_{n=0}^\infty\sigma_n\ee^{-a_n/h}f_n(h),
\]
其中 $\sigma_n$ 是常数（称为 Stokes 常数），$a_n$ 是 Borel 变换的奇点，$f_n$ 是形式幂级数。函数 $\tilde F$ 是 trans-级数（trans-series）的典型例子。单个 $f_n$ 可能不可 Borel 求和，但整个 trans-级数 $\tilde F$ 在许多情形下是 Borel 可和的。trans-级数的精确求和常常编码非微扰信息，尽管初始拟设是微扰的。一个漂亮而具体的例子——Airy 函数——在 [325, Sect. 3.1] 中有详细讨论。\end{remark}

\section{Schrödinger 方程的 Stokes 曲线}

现在我们限于 Schrödinger 方程，分析导致所谓 Stokes 曲线的精确 WKB 方法；Stokes 曲线结果是一个 $A_1$ 型谱网络。我们遵循 [216]。

考虑紧黎曼曲面 $S$，并选取典范丛的一个平方根 $K^{1/2}$。我们略微改变记号，考虑方程
\begin{equation}
(\partial^2-\eta^2t(z,\eta))\psi(z,\eta)=0,
\tag{16.21}
\end{equation}
其中 $\psi\in H^0(S,K^{-1/2})$，$t(z,\eta)=t_0(z)+\eta^{-1}t_1(z)+\dots$。我们通过 $\eta=h^{-1}$ 恢复上面的 Schrödinger 方程。

可以验证 $\varphi(z)=t_0(z)\mathrm{d}z^2$ 是 $S$ 上的亚纯二次微分。这个二次微分支配 WKB 方法的行为。注意

% [原文 p.194]
$\varphi$ 的极点是微分方程的奇点。$\varphi$ 的零点（称为转折点）也起重要作用。记 $P_0$ 与 $P_\infty$ 为 $\varphi$ 的零点集与极点集，并令 $P=P_0\cup P_\infty$。

最后，我们施加一些技术性条件，以简化对可能轨道的分析以及下文的连接公式 (16.25)。我们假设 $\varphi$ 至少有一个零点、至少有一个极点，所有零点都是单零点，所有极点的阶至少为 2。

如前，我们考虑 WKB 拟设 $\psi=\exp\left(\int^z s(z,\eta)\mathrm{d}z\right)$，它给出两个形式局部解 $\psi_\pm$。

我们可以如下得到形式整体解。定义 $s$ 的奇部与偶部为 $s_{\mathrm{odd}}=\frac12(s_+-s_-)$ 与 $s_{\mathrm{even}}=\frac12(s_++s_-)$。可以验证 $s_{\mathrm{even}}=-\frac12\partial(\ln s_{\mathrm{odd}})$。此外，$s_{\mathrm{odd}}\mathrm{d}z$ 在谱曲线 $\Sigma$ 上分歧点之外整体良定义 [229, Cor. 2.17]。因此
\begin{equation}
\psi_\pm(z,\eta)=\frac{1}{\sqrt{s_{\mathrm{odd}}}}\exp\left(\pm\int^z s_{\mathrm{odd}}\mathrm{d}z\right)
\tag{16.22}
\end{equation}
是形式整体解。确实，$\partial\psi_+=(s_{\mathrm{odd}}-\frac12\partial(\ln s_{\mathrm{odd}}))\psi_+=s\psi_+$，从而 $\partial^2\psi_+=(s^2+\partial s)\psi_+=\eta^2t\psi_+$。对 $\psi_-$ 有类似的计算。我们可以进一步把该表达式展开为
\begin{equation}
\psi_\pm(z,\eta)=\exp\left(\pm\int^z\sqrt{t_0(z)}\mathrm{d}z\right)\eta^{-1/2}\sum_{k=0}^\infty\eta^{-k}\psi_{\pm,k}(z).
\tag{16.23}
\end{equation}
该表达式的主项由
\begin{equation}
\psi_\pm(z,\eta)=t_0^{-1/4}\eta^{-1/2}\exp\left(\pm\int^z \eta\sqrt{t_0}+\frac{t_1}{2\sqrt{t_0}}\right)O(1+\eta^{-1}),
\tag{16.24}
\end{equation}
给出，它常被称为 Schrödinger 方程的 WKB 近似。它来自 $s_{\mathrm{odd}}=\sqrt{t_0}+\eta^{-1}\frac{t_1}{2\sqrt{t_0}}+O(\eta^{-2})$。

% [Tutor Note: 式 (16.24) 指数中的被积式按排版只写了 $\eta\sqrt{t_0}$，但由上面 $s_{\mathrm{odd}}$ 的展开，$\eta\int^z s_{\mathrm{odd}}\mathrm{d}z$ 的下一阶项 $\frac{t_1}{2\sqrt{t_0}}$ 也应落在积分号内（原文省略了外层括号）。核对方法：由 (16.23) 取首阶并对 $\eta$ 展开，指数中恰为 $\pm\int^z\left(\eta\sqrt{t_0}+\frac{t_1}{2\sqrt{t_0}}\right)\mathrm{d}z$；首阶以外的因子给出 $t_0^{-1/4}$ 与 $O(1+\eta^{-1})$。]

注意 $s_{\mathrm{odd}}\mathrm{d}z$ 在 $\varphi$ 的极点处不可积。为补救这一点，定义 $s_{\mathrm{odd}}$ 的正则部分为
\[
s^{\mathrm{reg}}_{\mathrm{odd}}(z,\eta)\mathrm{d}z:=\left(s_{\mathrm{odd}}(z)-\eta\sqrt{t_0(z)}\right)\mathrm{d}z.
\]
现在我们使用 Borel 求和。固定角度 $\theta\in[0,2\pi)$。用 (16.20) 定义 (16.23) 的带角 $\theta$ 的 Borel 可和性。这将依赖于点 $z\in S$。

正如我们在第 8 章所见，对亚纯微分 $\varphi$ 与角 $\theta$，可以联系 $S$ 上的一个叶状结构
\[
F_\theta:=\ker\operatorname{Re}\left(\ee^{\ii\theta}\sqrt{\varphi}\right).
\]
回忆这意味着：在每一点 $z\in S\setminus P_0$（它不是 $\varphi$ 的零点），存在唯一方向 $L(z)$，使得对所有 $v\in L(z)$，$\ee^{\ii\theta}\varphi(v)$ 为实数。

% [原文 p.195]
把叶 $L(z)$ 对 $z\in S$ 积分就得到叶状结构 $F_\theta$。换言之，$F_\theta$ 的一片叶 $\gamma$ 由如下事实刻画：$\int_\gamma\sqrt{\varphi}$ 的辐角为常数。

由第 8 章我们知道该叶状结构的各种轨道。回忆：若一片叶的一个端点落在 $P_0$（即 $\varphi$ 的零点），则称该叶为临界的。临界叶也称为 Stokes 曲线或 WKB 曲线。鞍点（saddle）连接 $P_0$ 中的两个点，而正则临界叶（也称分离叶）连接 $P_0$ 中一点与 $P_\infty$ 中一点。我们假设至多有一个鞍点（后文我们甚至假设没有鞍点）。

临界图或 Stokes 图是以 $P$ 为顶点集、以 Stokes 曲线为边的图。

Stokes 图的补集有三类区域 [66, Sect. 3.4]：带、半平面或圆盘；见图 16.2。这些恰是黎曼曲面上亚纯二次微分的临界轨道（参看第 8 章）。

\Fig[0.6\linewidth]{figures/fig_16_2.png}{16.2}{Stokes 图的补集。}

下一定理给出精确 WKB 方法与 Stokes 图（一个 $A_1$ 型谱网络）之间的重要联系：

\begin{theorem}[16.18（Koike–Schäfke）] 对 Stokes 图补集中的每个 $z\in S$，正则部分 $s^{\mathrm{reg}}_{\mathrm{odd}}(z,\eta)$ 在方向 $\theta$ 上 Borel 可和。其 Borel 和在 $\{\eta\in\C\mid|\arg(\eta)-\theta|<\pi/2,\ |\eta|\gg1\}$ 上以及在 $z$ 附近全纯。最后，若 $p\in P_\infty$ 是极点，则 $\int_p^z s^{\mathrm{reg}}_{\mathrm{odd}}(z,\eta)\mathrm{d}z$ 具有同样性质的 Borel 可和性。\end{theorem}

该定理由 T. Koike 在 2010 年 9 月 RIMS 研讨会“Exact WKB analysis – Borel summability of WKB solutions”上告知。作者们从未发表该结果。在二次位势 $t(z,\eta)=t_0(z)+\eta^{-1}t_1(z)+\eta^{-2}t_2(z)$ 的情形，这已在 [285] 中证明。一般情形最近由 Nikolaev [289, 290] 用几何方法证明。

最终我们要证明 $\psi_\pm$ 的 Borel 可和性。由于 Borel 可和性对乘积与指数保持（由命题 16.14），我们只需分析 $\int^z s^{\mathrm{reg}}_{\mathrm{odd}}(z,\eta)\mathrm{d}z$，其中我们沿谱曲线 $\Sigma$ 中的路径积分。

称 $\Sigma$ 中的一条路径为可容许路径，如果它到 $S$ 的投影只与 Stokes 图在极点处相交。

\begin{proposition}[16.19（{{[216, Lem. 2.20]}}）] 若路径 $\beta$ 的端点不投影到零点，且 $\beta$ 从不与鞍点相交，则 $\beta$ 是可容许路径之和。\end{proposition}

% [原文 p.196]
\begin{corollary}[16.20] 若 $\beta$ 是可容许路径之和，则 $\int_\beta s^{\mathrm{reg}}_{\mathrm{odd}}(z,\eta)\mathrm{d}z$ 是 Borel 可和的。\end{corollary}

特别地，若不存在鞍点轨道，则规范化 WKB 解是 Borel 可和的。

总结起来，固定角度 $\theta$ 给出一个 Stokes 图，它是与二次微分 $\varphi(z)=t_0(z)\mathrm{d}z^2$ 及角 $\theta$ 相联系的谱网络。在互补区域中，WKB 方法的形式解是 Borel 可和的，从而给出 Schrödinger 方程的精确解析解。

这一图景对高阶微分方程猜想上仍然成立：WKB 方法的形式解在谱网络的互补区域中是 Borel 可和的。不幸的是，高阶微分方程的分析要困难得多。在 opers 的特殊情形可以给出更精确的陈述，见下文第 16.8 节。

有两个有趣的现象值得研究：穿过一条 Stokes 曲线时会发生什么？变动角度 $\theta$ 时会发生什么？

\section{连接公式与 Voros 符号}

我们延续上一小节的设定。

假设 Stokes 图没有鞍点。则我们可以把 Borel 和在两个相邻区域 $D_1$ 与 $D_2$ 之间连接起来。记 $\Psi^{D_i}_\pm$ 为 $D_i$ 中的两个精确解（$i=1,2$）。记 $C$ 为 $D_1$ 与 $D_2$ 之间的 Stokes 曲线，见图 16.3。

\Fig[0.4\linewidth]{figures/fig_16_3.png}{16.3}{连接公式。}

不同区域中的解可由如下命题相互联系。这类公式称为连接公式。

\begin{proposition}[16.21（{{[216, Thm. 2.25]}}）] 若在 $C$ 上 $\operatorname{Re}\left(\ee^{\ii\theta}\int\sqrt{\varphi}\right)\ge0$，则
\begin{equation}
\Psi^{D_1}_+=\Psi^{D_2}_++\ii\Psi^{D_2}_-\qquad\text{和}\qquad \Psi^{D_1}_-=\Psi^{D_2}_-.
\tag{16.25}
\end{equation}
\end{proposition}

% [原文 p.197]
否则，我们得到：
\begin{equation}
\Psi^{D_1}_+=\Psi^{D_2}_+\qquad\text{和}\qquad \Psi^{D_1}_-=\Psi^{D_2}_-+\ii\Psi^{D_2}_+.
\tag{16.26}
\end{equation}

% [Tutor Note: (16.25)–(16.26) 中的系数 $\pm\ii$ 是 Stokes 常数在这一特殊情形（$\SL_2$、穿越单条 Stokes 曲线）的取值。把 (16.25) 写成矩阵形式即 $\begin{pmatrix}\Psi^{D_1}_+\\ \Psi^{D_1}_-\end{pmatrix}=\begin{pmatrix}1&\ii\\0&1\end{pmatrix}\begin{pmatrix}\Psi^{D_2}_+\\ \Psi^{D_2}_-\end{pmatrix}$，(16.26) 则为下三角 $\begin{pmatrix}1&0\\ \ii&1\end{pmatrix}$；这与 $\SL_2$ 的上/下三角 Stokes 因子形式一致，穿越方向相反时取互为转置的矩阵。]

下面分析变动角度 $\theta$ 时 Stokes 图与精确解的变化。

当至多有一个鞍点时，我们已经看到互补区域是带、半平面或圆盘。变动 $\theta$ 可能引入新的鞍点连接。Stokes 图的拓扑有两种可能的变化：flip 与 pop，见图 16.4 与图 16.5。

\Fig[0.75\linewidth]{figures/fig_16_4.png}{16.4}{Stokes 图的变化：flip。}

\Fig[0.65\linewidth]{figures/fig_16_5.png}{16.5}{Stokes 图的变化：pop。}

为描述精确 WKB 解如何变化，我们考虑称为 Voros 符号的特殊函数。记 $\tilde P$、$\tilde P_0$ 与 $\tilde P_\infty$ 为 $P=P_0\cup P_\infty$、$P_0$ 与 $P_\infty$ 在投影映射 $\pi:\Sigma\to S$ 下的原像。设 $\gamma\in H_1(\Sigma\setminus\tilde P;\Z)$，$\beta\in H_1(\Sigma\setminus\tilde P_0,\tilde P_\infty;\Z)$，即 $\Sigma\setminus\tilde P_0$ 中端点落在 $\tilde P_\infty$ 的路径。

联系于 $\beta$ 与 $\gamma$ 的 Voros 符号，或周期坐标，定义为
\[
W_\beta:=\exp\left(\int_\beta s^{\mathrm{reg}}_{\mathrm{odd}}\mathrm{d}z\right)
\qquad\text{和}\qquad
V_\gamma:=\exp\left(\int_\gamma s^{\mathrm{reg}}_{\mathrm{odd}}\mathrm{d}z\right).
\]
这里我们用精确 WKB 方法使 $W_\beta$ 与 $V_\gamma$ 有意义。注意 Voros 符号是良定义的，因为被积式 $s^{\mathrm{reg}}_{\mathrm{odd}}\mathrm{d}z$ 在 $\Sigma$ 上的对合（$\Sigma$ 是 $S$ 的分歧二重覆盖）下不变。尤其重要的是 $\gamma$ 为绕鞍点轨道的闭链时的 Voros 符号 $V_\gamma$；见图 16.6。

% [原文 p.198]
\Fig[0.85\linewidth]{figures/fig_16_6.png}{16.6}{Voros 符号。}

Voros 符号是平坦联络空间上的坐标。在 flip 之下，它们按 X-变异规则（第 5 章的式 (5.1)）变化。因此它们可以等同于 $\SL_2(\C)$-特征簇的 Fock–Goncharov 坐标；见第 10.2 节。这是 Iwaki–Nakanishi [216] 的主要结果。

这一结论对更高秩猜想上仍然成立：已知 Fock–Goncharov 坐标可由谱曲线的周期描述。该谱曲线应当与 WKB 方法中出现的谱曲线相同，而特殊闭曲线附近的精确 WKB 解应当等同于 Fock–Goncharov 坐标。

\section{Opers}

黎曼曲面上微分方程的一个重要特殊例子由 opers 给出。它们把 Schrödinger 方程推广到高阶微分方程，并允许表述精确的猜想。

再次从定义在开集 $U\subset\C$ 上的单变量 $z$ 的 $n$ 阶微分方程出发：
\[
(a_0(z,h)\partial^n+a_1(z,h)\partial^{n-1}+\dots+a_{n-1}(z,h)\partial+a_n(z,h))\psi(z,h)=0,
\]
其中系数为全纯函数（也可考虑更一般的函数），且 $a_0$ 处处非零。全纯函数空间通过前复合与后复合作用在这些算子组成的空间上。利用这两个作用可以把 $a_0$ 化为 1（除以 $1/a_0$），并把 $a_1$ 化为 0（通过适当的平移）。

其余参数 $(a_2,a_3,\dots,a_n)$ 在把坐标 $z$ 变为 $w(z)$ 时并不像张量那样变换。它们的变换行为是复杂的，但有一种方式把这些系数及其导数组合起来得到张量。为此，再次考虑微分算子的矩阵版本（如式 (16.10)）：

% [原文 p.199]
\begin{equation}
\partial+
\begin{pmatrix}
0&1&0&0\\
\vdots&\ddots&\ddots&0\\
0&\cdots&0&1\\
-a_n&-a_{n-1}&\cdots&-a_1
\end{pmatrix}.
\tag{16.27}
\end{equation}

\begin{proposition}[16.22（{{[107, Thm. 1]}}）] 存在一个幂幺规范变换，把 (16.27) 中的联络变为如下形式的联络：
\begin{equation}
\partial+
\begin{pmatrix}
0&1&0&\cdots&0\\
-u_2&0&1&\ddots&\vdots\\
-u_3&-u_2&0&\ddots&0\\
\vdots&&\ddots&\ddots&1\\
-u_n&\cdots&-u_3&-u_2&0
\end{pmatrix}.
\tag{16.28}
\end{equation}
\end{proposition}

上面的 $u_i$（$i=2,\dots,n$）是 $(a_j)_{j\ge3}$ 及其导数的线性组合，以及 $a_2$ 及其导数的多项式。对 $i\ge3$，坐标 $u_i$ 是 $(i,0)$ 型张量，即它是一个 $i$-微分。坐标 $u_2$ 像复射影结构那样变换：
\begin{equation}
u_2(w)=\left(\frac{\mathrm{d}z}{\mathrm{d}w}\right)^2u_2(z)+S(z,w),
\tag{16.29}
\end{equation}
其中 $S(z,w)$ 表示 Schwarz 导数。

% [Tutor Note: 为何 $u_2$ 与 $u_{i\ge3}$ 要分开处理：$i\ge3$ 时 $u_i$ 是 $i$-微分，在坐标变换下按 $u_i(w)=u_i(z)(\mathrm{d}z/\mathrm{d}w)^i$ 张量式地变换；而 $u_2$ 的变换 (16.29) 多出一个 Schwarzian 修正项，故它不是 2-微分，而应视为复射影结构（一个“联络”）。这也解释了命题 16.22 的措辞：$u_3,\dots,u_n$ 给出微分，而 $u_2$ 单独编码 $\C\PP^1$-结构。]

我们邀请读者对 $n=2$ 做计算：算出 Schrödinger 算子 $\partial_z^2+a_2(z)$ 在坐标变换 $z\mapsto w(z)$ 下如何变化。必须使用全纯函数的双重作用，把所得算子再次化为 Schrödinger 算子的形式。于是 $a_2$ 的变换规则 (16.29) 出现。Schwarz 导数
\[
S(f,z)=\frac{f'''}{f'}-\frac32\left(\frac{f''}{f'}\right)^2
\]
是在 Möbius 变换下不变的微分算子。关于 Schwarz 导数的细节，我们推荐 [293]。

命题 16.22 用一族次数从 3 到 $n$ 的微分，以及一个编码复射影结构的特殊参数，来参数化微分方程的空间。

曲面 $S$ 上的复射影结构（简称 $\C\PP^1$-结构）是一个图册，其坐标卡映入 $\C\PP^1$ 的开子集，转移函数为 Möbius 变换。复射影结构诱导一个复结构（因为 Möbius 变换全纯）。已知所有

% [原文 p.200]
诱导 $S$ 上同一复结构的 $\C\PP^1$-结构构成一个仿射空间，其模型向量空间是全纯二次微分空间 $H^0(C,K_C^{\otimes2})$。

因此，若从黎曼曲面 $C$ 出发并选取一个与 $C$ 上复结构相容的复射影结构，则任何其他相容的 $\C\PP^1$-结构都可由一个全纯二次微分描述。于是，在带 $\C\PP^1$-结构的黎曼曲面上，$n$ 阶微分方程可由全纯微分的一个元组 $(u_2,u_3,\dots,u_n)$ 参数化，其中 $u_i\in H^0(C,K^{\otimes i})$。

我们可以反解命题 16.22 给出的幂幺规范变换，用全纯微分写出微分方程。对小的 $n$，我们得到（另见 [107, Tab. 1]）：
\[
\begin{aligned}
n=2&:\quad \partial^2+u_2\\
n=3&:\quad \partial^3+u_2\partial+u_3+\tfrac12\partial u_2\\
n=4&:\quad \partial^4+u_2\partial^2+(u_3+\partial u_2)\partial+\left(u_4+\tfrac12\partial u_3+\tfrac1{20}\partial^2u_2+\tfrac9{100}u_2^2\right)
\end{aligned}
\]
黎曼曲面 $C$ 上这种形式的微分算子称为 opers。它们由 Beilinson 与 Drinfeld 在 [44] 中引入。正如我们在例 16.3 之后所见，这些微分算子作用在空间 $K^{(1-n)/2}$ 上。对偶数 $n$，这意味着我们还须固定典范丛的一个平方根，即 $C$ 上的一个自旋结构。

总结起来，我们考虑黎曼曲面 $C$，连同相容的 $\C\PP^1$-结构与自旋结构。对 $u=(u_2,\dots,u_n)$（$C$ 上次数从 2 到 $n$ 的全纯微分元组），记 $D_u:K^{(1-n)/2}\to K^{(n-1)/2}$ 为以 $u$ 参数化的 oper。

为进行 WKB 分析，我们引入参数 $h\in\C^*$，按
\[
h.u=(h^{-2}u_2,h^{-3}u_3,\dots,h^{-n}u_n)
\]
重标度各微分，并考虑 oper $D_{h.u}$。乘以 $h$ 的某次幂后，这相当于把 $D_u$ 中所有导数 $\partial$ 换成 $h\partial$。于是我们得到一个 $h$-联络。

% [Tutor Note: 核对缩放 $h.u$：由 $D_{h.u}=\partial^n+\sum_{i=2}^{n}u_ih^{-i}\partial^{n-i}$，乘以 $h^n$ 得 $h^nD_{h.u}=h^n\partial^n+\sum_{i=2}^{n}u_ih^{n-i}\partial^{n-i}$，这正是把 $D_u=\partial^n+\sum_{i=2}^{n}u_i\partial^{n-i}$ 中所有 $\partial$ 换为 $h\partial$ 的结果。故 $D_{h.u}$ 与 $D_u$ 只差一个整体因子，而 $\partial\mapsto h\partial$ 使 $D_{h.u}$ 成为 $h$-联络（各矩阵元与 $h\partial$ 同阶）。特例 $n=2$：$h^2D_{h.u}=h^2\partial^2+u_2$。]

现在我们可以对 opers 情形下精确 WKB 方法的收敛性表述一个精确的猜想。

\begin{conjecture}[16.23] 设 $C$ 是带相容 $\C\PP^1$-结构与自旋结构的黎曼曲面。对微分元组 $u=(u_2,\dots,u_n)\in\bigoplus_{i=2}^{n}H^0(C,K^{\otimes i})$，联系带角 $\theta$ 的谱网络 $\mathcal{W}^\theta_u$ 与如上 oper $D_{h.u}$。
\end{conjecture}

则把方向 $\theta$ 上的精确 WKB 方法应用于 $D_{h.u}$ 所对应的微分方程，该过程在 $\mathcal{W}^\theta_u$ 的补集上收敛。此外，对 $\gamma\in H_1(\Sigma)$，谱曲线 $\Sigma$ 上阿贝尔化局部系统的和乐 $X^\theta_\gamma(h)\in\C^*$ 在 $h\to0$ 时渐近地表现为
\begin{equation}
X^\theta_\gamma(h)\sim\exp(Z_\gamma/h),
\tag{16.30}
\end{equation}
其中 $Z_\gamma=\int_\gamma\lambda\in\C$ 是 Liouville 形式 $\lambda$ 绕 $\gamma$ 的周期。

\begin{remark}[16.24] Hitchin 截面是把平坦联络联系于如上微分元组 $u$ 的另一种方式。寻找平坦截面再次给出一个微分

% [原文 p.201]
方程，它由 $u$ 参数化。这个微分方程不是显式的，因为 Hitchin 截面使用了高度超越的非阿贝尔 Hodge 对应。尽管如此，可以预期上述猜想的陈述对 Hitchin 截面仍成立，且 (16.30) 的渐近行为是指数好的。部分结果例如可在 [280] 中找到。\end{remark}

% [原文 p.202]
% 本页（印刷 p.202）为章末空白页，原文无正文，故此处不置内容。
